\section{The Connection Between the Net Change Theorem, Differential Equations, and Euler's Method}

The connection between the net change theorem, differential equations, and Euler's method is structural. They are different views of the same fundamental operation:

\begin{itemize}
  \item Differential equations describe \emph{local rates of change}.
  \item Integration accumulates those local changes into a \emph{global change}.
  \item Euler's method numerically performs that accumulation step-by-step.
\end{itemize}

The net change theorem is the bridge connecting these ideas.

Suppose that
\[
  F'(t)=f(t).
\]

Then the net change theorem states:
\[
  F(b)-F(a)=\int_{a}^{b} f(t)\,dt \quad \text{or}\quad F(b)=F(a)+\int_{a}^{b} f(t)\,dt.
\]

The interpretation is:

\begin{itemize}
  \item $f(t)$ is the instantaneous rate of change,
  \item the integral accumulates infinitesimal changes,
  \item the result is the total change in $F$.
\end{itemize}

Now consider a differential equation:
\[
  y'(t)=f(t,y(t)).
\]

This equation states that the instantaneous change of the unknown function $y$ is determined by the function $f$.

Applying the net change theorem to the unknown solution $y$, we integrate both sides from $t_{0}$ to $t$:
\[
  y(t)-y(t_{0})=\int_{t_0}^{t} f(s,y(s))\,ds.
\]

If the initial value is
\[
  y(t_{0})=y_{0},
\]
then:
\[
  y(t)=y_{0}+\int_{t_0}^{t} f(s,y(s))\,ds.
\]

This equation has a direct interpretation:

\begin{itemize}
  \item Start with the initial value $y_{0}$,
  \item accumulate all infinitesimal changes $f(s,y(s))\,ds$,
  \item obtain the final value $y(t)$.
\end{itemize}

Thus, a differential equation is fundamentally an accumulation law.

\subsection{Euler's Method as Numerical Integration}

Euler's method arises by approximating the integral with a Riemann sum. Take a small step size $h$. Over the interval $[t_{n},t_{n}+h]$,
\[
  y(t_{n+1})-y(t_{n}) = \int_{t_n}^{t_n+h}f(s,y(s))\,ds.
\]

Euler's approximation assumes that the slope remains approximately constant over this short interval:
\[
  f(s,y(s)) \approx f(t_{n},y_{n}).
\]

Therefore,
\[
  \int_{t_n}^{t_n+h}f(s,y(s))\,ds \approx h\,f(t_{n},y_{n}).
\]

Substituting this approximation gives Euler's update rule:
\[
  y_{n+1}=y_{n}+h\,f(t_{n},y_{n}).
\]

This is not merely inspired by integration. It \emph{is} numerical integration. Each Euler step performs the following process:

\begin{enumerate}
  \item Measure the current rate of change,
  \item assume it remains roughly constant over a short interval,
  \item accumulate the resulting small change,
  \item repeat.
\end{enumerate}

Thus, every Euler step is a discrete version of the net change theorem:
\[
  \text{new value}= \text{old value}+ \text{accumulated change}.
\]

This explains why:

\begin{itemize}
  \item numerical ODE solving is often called \emph{integration},
  \item solution curves are called \emph{integral curves},
  \item ODE solvers are often called \emph{integrators}.
\end{itemize}

